\documentclass[11pt,a4paper]{article}

\usepackage[utf8]{inputenc}
\usepackage[T1]{fontenc}
\usepackage[provide=*,english]{babel}
\usepackage{lmodern}
\usepackage{microtype}
\usepackage{amsmath,amssymb,mathtools}
\usepackage{geometry}
\usepackage{xcolor}
\usepackage{enumitem}
\usepackage{booktabs}
\usepackage{fancyhdr}
\usepackage[hidelinks]{hyperref}

\geometry{left=2.3cm,right=2.3cm,top=2.2cm,bottom=2.25cm,headheight=14pt}
\setlength{\parindent}{0pt}
\setlength{\parskip}{0.38em}
\setlist[itemize]{topsep=0.25em,itemsep=0.18em,leftmargin=1.5em}
\allowdisplaybreaks

\definecolor{blue}{RGB}{34,73,110}
\definecolor{gray}{RGB}{90,96,102}

\pagestyle{fancy}
\fancyhf{}
\fancyhead[L]{\small\color{gray}GC2D: ABBA method}
\fancyhead[R]{\small\color{gray}ABBA2 implicit: simultaneous formulation}
\fancyfoot[C]{\thepage}
\renewcommand{\headrulewidth}{0.3pt}

\newcommand{\R}{\mathbb{R}}
\newcommand{\Id}{I}
\newcommand{\Y}{\mathcal{Y}}
\newcommand{\M}{\mathcal{M}}
\newcommand{\Om}{\Omega_{\times}}
\newcommand{\PhiAB}{\Phi^{\mathrm{ABBA}}_{h,t_n}}
\newcommand{\norm}[1]{\left\lVert #1\right\rVert}

\begin{document}

\begin{center}
  {\LARGE\bfseries\color{blue}ABBA2 Implicit: Simultaneous State--Multiplier Formulation}\par
  \vspace{0.35em}
  {\large Simultaneous Newton formulation for the output state and multiplier}\par
  \vspace{0.45em}
  {\small GC2D implementation of equation (21) from the original implicit note}
\end{center}

\vspace{0.5em}

\section{Purpose and relation to the reduced formulation}

The \texttt{simultaneous\_state\_multiplier} option of
\texttt{ABBA2Implicit} computes the same exact projected map as its
\texttt{reduced\_multiplier} option but uses a different nonlinear system. The
reduced formulation eliminates the duplicated output state and solves only for
the two-component multiplier. The simultaneous formulation retains both the
four-component duplicated output and the multiplier as Newton
unknowns. Consequently, for one guiding-centre particle, each Newton correction
solves a $6\times6$ system rather than the reduced $2\times2$ Schur complement.

The two formulations must converge to the same physical state whenever the
projection root is locally unique and both nonlinear solves converge to that root.
Finite tolerances and floating-point linear algebra can produce differences at the
round-off level.

\section{Guiding-centre and duplicated states}

The physical state $z=(X,Y)\in\R^2$ satisfies
\[
 \dot z=f(z,t),
 \qquad
 f(z,t)=J_0\nabla_z\phi(z,t),
 \qquad
 J_0=
 \begin{pmatrix}
  0&-1\\
  1&0
 \end{pmatrix}.
\]
Its spatial Jacobian
\[
 W(z,t)=D_zf(z,t)=J_0\nabla_z^2\phi(z,t)
\]
satisfies $W^TJ_0+J_0W=0$ because the Hessian of $\phi$ is symmetric.

The duplicated state and physical diagonal are
\[
 \Y=(u,v)\in\R^4,
 \qquad
 \M=\{(u,v):u=v\}.
\]
Define
\[
 G=(\Id_2\;-\Id_2),
 \qquad
 N=G^T=
 \begin{pmatrix}\Id_2\\-\Id_2\end{pmatrix},
 \qquad
 E=
 \begin{pmatrix}\Id_2\\\Id_2\end{pmatrix}.
\]
Then $Ez=(z,z)$ is diagonal, $N\mu=(\mu,-\mu)$ is a correction
normal to the diagonal, and
\[
 GE=0,
 \qquad
 GN=GG^T=2\Id_2.
\]

\section{Endpoint-time ABBA map}

Let $s=h/2$ and $t_{n+1}=t_n+h$. The nonautonomous ABBA map is
\[
 \PhiAB
 =A_{s,t_{n+1}}\circ B_{s,t_{n+1}}
  \circ B_{s,t_n}\circ A_{s,t_n},
\]
where
\[
 A_{s,t}(u,v)=\bigl(u+s f(v,t),v\bigr),
 \qquad
 B_{s,t}(u,v)=\bigl(u,v+s f(u,t)\bigr).
\]
Starting from $(u^{(0)},v^{(0)})$, the four explicit stages are
\begin{align*}
 u^{(1)}&=u^{(0)}+s f(v^{(0)},t_n),\\
 v^{(1)}&=v^{(0)}+s f(u^{(1)},t_n),\\
 v^{(2)}&=v^{(1)}+s f(u^{(1)},t_{n+1}),\\
 u^{(2)}&=u^{(1)}+s f(v^{(2)},t_{n+1}).
\end{align*}
Thus $\PhiAB(u^{(0)},v^{(0)})=(u^{(2)},v^{(2)})$.
The endpoint-time order makes the map symmetric:
\[
 \Phi^{\mathrm{ABBA}}_{-h,t_{n+1}}
 =\left(\Phi^{\mathrm{ABBA}}_{h,t_n}\right)^{-1}.
\]

At the stage points, define
\[
 W_1=W(v^{(0)},t_n),\quad
 W_2=W(u^{(1)},t_n),\quad
 W_3=W(u^{(1)},t_{n+1}),\quad
 W_4=W(v^{(2)},t_{n+1}),
\]
and $S=W_2+W_3$. The exact duplicated ABBA Jacobian is
\[
 D\PhiAB=
 \begin{pmatrix}
  \Id_2+s^2W_4S
  &s(W_1+W_4)+s^3W_4SW_1\\[2mm]
  sS
  &\Id_2+s^2SW_1
 \end{pmatrix}.
\]
The matrix products retain this order because stage Jacobians evaluated at
different points need not commute.

\section{Simultaneous projection equations}

Given the diagonal input $\Y_n=Ez_n$, pre-correct the two copies by the
unknown multiplier:
\[
 \widehat{\Y}_n(\mu)=\Y_n+N\mu=(z_n+\mu,z_n-\mu).
\]
The simultaneous formulation retains the diagonal output $\Y_{n+1}$ as an additional
unknown and imposes
\[
 d(\Y_{n+1},\mu)
 :=\Y_{n+1}-N\mu
   -\PhiAB(\Y_n+N\mu)=0,
\]
\[
 g(\Y_{n+1}):=G\Y_{n+1}=0.
\]
The first equation states that the same multiplier used to separate the input
copies is applied with the opposite placement around the ABBA output. The second
equation constrains the final copies to the physical diagonal.

Writing $\Y_{n+1}=(y_u,y_v)$, the residual components used by the implementation are
\[
 d_u=y_u-\mu-u^{(2)}(\mu),
 \qquad
 d_v=y_v+\mu-v^{(2)}(\mu),
 \qquad
 g=y_u-y_v.
\]

\section{Exact simultaneous Newton system}

Because $G$ and $N=G^T$ are constant,
\[
 D_{\Y_{n+1}}d=\Id_4,
 \qquad
 D_\mu d=-\left[\Id_4+D\PhiAB(\widehat{\Y}_n)\right]N,
\]
\[
 D_{\Y_{n+1}}g=G,
 \qquad
 D_\mu g=0.
\]
Therefore, at Newton iterate $(\Y_{n+1}^{(k)},\mu^{(k)})$, the simultaneous formulation
solves
\[
 \boxed{
 \begin{pmatrix}
  \Id_4&-\left[\Id_4+D\PhiAB(\widehat{\Y}_n)\right]G^T\\[2pt]
  G&0
 \end{pmatrix}
 \begin{pmatrix}
  \Delta\Y_{n+1}\\
  \Delta\mu
 \end{pmatrix}
 =-
 \begin{pmatrix}
  d(\Y_{n+1},\mu)\\
  g(\Y_{n+1})
 \end{pmatrix}.
 } \tag{21}
\]
The number (21) is retained deliberately to identify the equation from the
original \texttt{ABBA\_implicit} document.

The update is
\[
 \Y_{n+1}^{(k+1)}=\Y_{n+1}^{(k)}+\Delta\Y_{n+1},
 \qquad
 \mu^{(k+1)}=\mu^{(k)}+\Delta\mu.
\]
No inverse is formed explicitly; a dense linear solve is used for each
$6\times6$ particle block.

\section{Schur complement and equivalence to the reduced formulation}

Eliminating $\Delta\Y_{n+1}$ from (21) gives the Schur complement
\[
 \mathcal K
 =G\left[\Id_4+D\PhiAB(\widehat{\Y}_n)\right]G^T
 =2\Id_2+GD\PhiAB(\widehat{\Y}_n)G^T.
\]
Let the reduced residual be
\[
 R(\mu)=G\PhiAB(\Y_n+G^T\mu)+2\mu.
\]
Then
\[
 \mathcal K=R'(\mu).
\]
Moreover, from the definitions of $d$ and $g$,
\[
 Gd=g-R(\mu).
\]
The second block row of (21), after substituting the first one, therefore yields
\[
 \mathcal K\,\Delta\mu=-R(\mu).
\]
This is exactly the Newton correction used by the reduced formulation. Hence both
algorithms generate the same multiplier iterates in exact arithmetic when they start
from the same $\mu^{(0)}$, even though the simultaneous formulation also updates the duplicated
output unknown.

\section{Initial iterate and convergence criterion}

The implementation uses
\[
 \mu^{(0)}=0,
 \qquad
 \Y_{n+1}^{(0)}=\PhiAB(\Y_n).
\]
This makes $d=0$ initially while leaving the diagonal defect $g$ to drive the first
Newton correction. For a physical input scale
\[
 s_z=\max\bigl(1,\norm{z_n}_\infty\bigr),
\]
the iteration stops when
\[
 \max\left(\norm{d}_\infty,\norm{g}_\infty\right)
 \leq \text{atol}+\text{rtol}\,s_z.
\]
Newton is not damped, so the residual norm is not required to decrease monotonically.
Failure to converge within the configured iteration count, or a singular simultaneous
Jacobian, terminates the integration step with an error.

For $N_p$ independent guiding-centre particles, physical states use component-major
storage,
\[
 (X_1,\ldots,X_{N_p},Y_1,\ldots,Y_{N_p}),
\]
but equation (21) is assembled as $N_p$ independent $6\times6$ systems. This preserves
the absence of artificial cross-particle derivatives while allowing a vectorized solve.

At convergence, $y_u$ and $y_v$ agree up to the stopping tolerance. The physical state is
represented neutrally by
\[
 z_{n+1}=\frac{y_u+y_v}{2}.
\]

\section{Symmetry and symplecticity of the converged map}

At the exact root, both formulations satisfy
\[
 \Y_{n+1}-N\mu
 =\Phi^{\mathrm{ABBA}}_{h,t_n}(\Y_n+N\mu),
 \qquad
 \Y_n,\Y_{n+1}\in\M.
\]
Using the inverse ABBA map and defining the reverse multiplier as
$\mu_{\mathrm{inv}}=-\mu$ gives
\[
 \Psi_{-h,t_{n+1}}\circ\Psi_{h,t_n}=\Id,
\]
provided the root is locally unique and solved exactly.

The duplicated symplectic form is
\[
 \Om=
 \begin{pmatrix}
  0&J_0\\
  J_0&0
 \end{pmatrix}.
\]
Every $A$ and $B$ shear preserves $\Om$, and therefore so does $\PhiAB$. The
diagonal and correction directions satisfy
\[
 E^T\Om E=2J_0,
 \qquad
 N^T\Om N=-2J_0,
 \qquad
 E^T\Om N=N^T\Om E=0.
\]
Differentiating the exact projected equation with
$M=D\mu_h(z)$ and $J_\Psi=D\Psi_{h,t_n}(z)$ gives
\[
 D\PhiAB(E+NM)=EJ_\Psi-NM.
\]
Preservation of $\Om$ and the orthogonality identities cancel the terms containing
$M$, leaving
\[
 \boxed{J_\Psi^TJ_0J_\Psi=J_0.}
\]
Thus the simultaneous formulation is the same symmetric and symplectic physical
map as the reduced formulation at the exact root. In finite arithmetic, measured reversibility and symplecticity
also include the simultaneous Newton tolerance and linear-solve round-off.

\section{Implementation correspondence}

The production implementation follows the mathematical blocks directly:
\begin{itemize}
 \item the endpoint-time stages evaluate $\PhiAB$;
 \item their exact derivatives assemble $D\PhiAB$;
 \item each particle receives the $6\times6$ matrix in equation (21);
 \item diagnostics record the simultaneous residual norm, iteration count, and
       multiplier norm;
 \item the optional ideal physical tangent is evaluated by implicit differentiation
       of the common exact root, independently of the finite Newton stopping decision.
\end{itemize}

\end{document}
